\ifx\allfiles\undefined
%生成目录超链接，使用UTF8字符集
\documentclass[12pt,UTF8]{ctexart}


\usepackage{mathtools}%花括号括起的公式
\usepackage{enumerate}%列表
\usepackage{graphicx}%表格
\usepackage[table]{xcolor}%彩色表格
\usepackage{booktabs}%调整表格行高


\renewcommand\tabcolsep{50pt}
\renewcommand\arraystretch{2}

%文本框
\usepackage{tcolorbox}
\tcbuselibrary{skins,vignette,raster,listings,listingsutf8,theorems, breakable,magazine,fitting}

%页面设置
\usepackage{geometry}%设置页面
\usepackage{setspace}%设置局部行距

\usepackage{ctex}
\usepackage{CJK}

%图像处理
\usepackage{graphicx}
\usepackage{float}%允许浮动
\usepackage{caption}%设置标题
\usepackage{graphicx, subfig}

%附录
\usepackage[toc,page]{appendix}

%导言区
\title{中小微企业的信贷决策}

%不要页眉
\pagestyle{plain}
\bibliographystyle{plain}

\geometry{a4paper}%将纸张设置为A4大小

%修改文档结构层次的名称
\usepackage{titlesec}%chapter1修改为一
\titleformat{\part}{\centering\Huge\bfseries}{\thepart}{1em}{}
\renewcommand\thepart{\chinese{part}}

%subsubsubsection的使用
%\usepackage{titlesec}
%\titleclass{\subsubsubsection}{straight}[\subsection]
%
%\newcounter{subsubsubsection}[subsubsection]
%\renewcommand\thesubsubsubsection{\thesubsubsection.\arabic{subsubsubsection}}
%\renewcommand\theparagraph{\thesubsubsubsection.\arabic{paragraph}}
%
%\titleformat{\subsubsubsection}
%{\normalfont\normalsize\bfseries}{\thesubsubsubsection}{1em}{}
%\titlespacing*{\subsubsubsection}
%{0pt}{3.25ex plus 1ex minus .2ex}{1.5ex plus .2ex}
%
%\makeatletter
%\renewcommand\paragraph{\@startsection{paragraph}{5}{\z@}%
%	{3.25ex \@plus1ex \@minus.2ex}%
%	{-1em}%
%	{\normalfont\normalsize\bfseries}}
%\renewcommand\subparagraph{\@startsection{subparagraph}{6}{\parindent}%
%	{3.25ex \@plus1ex \@minus .2ex}%
%	{-1em}%
%	{\normalfont\normalsize\bfseries}}
%\def\toclevel@subsubsubsection{4}
%\def\toclevel@paragraph{5}
%\def\toclevel@paragraph{6}
%\def\l@subsubsubsection{\@dottedtocline{4}{7em}{4em}}
%\def\l@paragraph{\@dottedtocline{5}{10em}{5em}}
%\def\l@subparagraph{\@dottedtocline{6}{14em}{6em}}
%\makeatother
%
%\setcounter{secnumdepth}{4}

\begin{document}
\else
\fi
%分文件开始


%参考文献
\begin{thebibliography}{99}
	%在没有原样抄写的情况下，网站中的“#”、“_”改成“\#”、“\_”，前面加一个反斜杠
	%现已有原样抄写的环境下，可以直接将链接放入

	%开始：参考文献
	\bibitem{ref1} \verb|司守奎、孙兆亮，《数学建模算法与应用》（第二版）[M]，国防工业出版社，2016。|
	\bibitem{ref2} \verb|姜启源，《数学模型》[M]，高等教育出版社，1993。|
	%\bibitem{ref3} \verb|曹明，预测评价指标RMSE、MSE、MAE、MAPE、SMAPE，https://www.cnblogs.com/think90/articles/11786016.html，2020年9月12日。|
	%\bibitem{ref4} \verb|easyAI，决策树 – Decision tree，https://easyai.tech/ai-definition/decision-tree/，2020年09月12日。|
	\bibitem{ref3} \verb|新浪财经，新冠肺炎疫情对若干行业的影响分析[OL]，https://finance.sina.com.cn/review/jcgc/2020-02-08/doc-iimxyqvz1258307.shtml，2020年9月12日。|

	%\bibitem{ref4} \verb||
	%\bibitem{ref5} \verb||
	%\bibitem{ref6} \verb||

	%结束：参考文献
\end{thebibliography}
\addcontentsline{toc}{section}{参考文献}%将“参考文献”字样加入目录中


\newpage


%分文件结束
\ifx\allfiles\undefined
\end{document}
\fi